\section{A rowwise tensor model for the Thom--Whitney end}

\paragraph{Archive scope.}
These results concern cote 136, physical pp.~117, 119 and 121,
archive footers 116, 118 and 120, supplied scans 117, 119 and 121.
They use the inherited standard model AQ-136-DP-MODEL and the
preceding checked ARCH-136-TW-VERTICAL and ARCH-136-TW-TOR.
The integral common-source and localization questions remain open.

\paragraph{ARCH-136-TW-TENSOR.}
Let \(k\) be a commutative ring, \(S=\Gamma_k(T)\), and let \(C\) be a simplicial, nonnegatively weighted, nonnegative cochain divided-power \(S\)-algebra. Write \(R=GU\) for the right-cofree monad in the original simplicial direction, and
\[
 \mathcal T(C)_n=\int_{[r]\in\Delta} B_r\otimes_S(R^{r+1}C)_n .
\]
There is a natural strict divided-power map
\[
 \lambda_C:B\otimes_S C\longrightarrow\mathcal T(C)
\]
which is a quasi-isomorphism in the original simplicial direction in every fixed weight and form degree. At each original simplex and weight it is also a homotopy equivalence of vertical cochain complexes.

\begin{proof}
Use the first-vertex map of the linear-end lemma proved in ARCH-136-TW-TOR. At original simplex \(n\), a resolution string
\([l_r]\xrightarrow{\alpha_r}\cdots\xrightarrow{\alpha_1}[l_0]\xrightarrow{\theta}[n]\)
determines \(\beta:[r]\to[n]\) by
\(\beta(i)=\theta\alpha_1\cdots\alpha_i(0)\).
Its coordinate map is
\[
 b\otimes c\longmapsto B(\beta)b\otimes
 C(\theta\alpha_1\cdots\alpha_r)c .
\]
Both factors are strict divided-power algebra maps, so their tensor is one. The string face and degeneracy identities established in the linear-end lemma prove the end relations, simplicial naturality, and compatibility with the cochain differential.

For a fixed form degree \(i\), set \(P=B^{*,i}\). In each simplicial degree it is the free graded \(S\)-module with basis
\(\gamma_{a_1}(x_1)\cdots\gamma_{a_r}(x_r)e_I\), \(|I|=i\).
For a weight cutoff \(w\), retain the basis elements of relative polynomial degree at most \(w-i\), taking the zero module if \(w<i\). This is a simplicial submodule and agrees with \(P\) in every total weight at most \(w\). Filter it by relative polynomial degree \(b\). The associated graded is
\[
 S\otimes_k\Gamma_k^b(Q)\otimes_k\Lambda_k^i(Q),
 \qquad Q=\operatorname{DK}(k[1]).
\]
In Dold--Kan coordinates a monomial in this module uses at most \(b+i\) coordinate positions. If the simplicial degree exceeds \(b+i\), one coordinate is omitted and the monomial is degenerate. Consequently normalization vanishes above \(b+i\), and hence above \(w\) for the finite filtration. The normalized terms are finite projective graded \(S\)-modules, being direct summands of finite free simplicial terms.

The linear-end lemma applies to this truncation with the simplicial module \(C^{*,j}\). Nonnegative weights make truncation harmless in total weight \(w\). Summing over the finitely many pairs \(i+j=q\) gives the claimed horizontal quasi-isomorphism in form degree \(q\).

The inclusion \(C_n\to B_n\otimes_S C_n\) is a vertical homotopy equivalence: use the usual \(S\)-linear ordered-coordinate contraction of \(B_n\) to its vertex \(S\), tensored with \(C_n\). The inclusion \(C_n\to\mathcal T(C)_n\) is the vertical homotopy equivalence of ARCH-136-TW-VERTICAL. Their triangle commutes because the displayed coordinate formula sends \(1\otimes c\) to the coaugmentation of \(c\). The two-out-of-three property in the homotopy category proves the vertical assertion.
\end{proof}

\paragraph{ARCH-136-TW-RATIONAL.}
Suppose \(k\) is a \(\mathbb Q\)-algebra and \(C\) is admissible. If
\(N(C^{u,j})\) is acyclic for \(u\geq\mu_C(j)\), then both
\(B\otimes_S C\) and \(\mathcal T(C)\) have acyclic horizontal form-\(q\) rows for
\[
 w\geq \max_{0\leq i\leq q}\{\mu_C(q-i)+i+1\}.
\]
They retain the vertical admissibility bounds of \(C\). In particular
\(C\to B\otimes_S C\leftarrow B\) is a natural strict divided-power common-target diagram of admissible theories.

\begin{proof}
For every \(i\), the normalized simplicial \(S\)-module \(NB^{*,i}\) is a projective resolution of \(k\), shifted by \(i\) in both homological degree and weight. Indeed the standard row calculation is zero in weights greater than \(i\), while in weight \(i\) the row is \(\Lambda^iQ\), whose normalization is \(k[i]\); weights below \(i\) vanish. Its terms are projective over \(S\), and its terms below simplicial degree \(i\) vanish.

Now \(S=\Gamma_k(T)\cong k[T]\), and the augmentation module \(k\) has the graded free resolution
\[
 0\longrightarrow S(-1)\xrightarrow{T}S\longrightarrow k\longrightarrow0.
\]
Eilenberg--Zilber and this resolution identify the contribution of
\(B^{*,i}\otimes_S C^{*,j}\) in total weight \(w\), up to the homological shift \(i\), with the cone of
\[
 T:N(C^{w-i-1,j})\longrightarrow N(C^{w-i,j}).
\]
Both complexes are acyclic whenever \(w-i-1\geq\mu_C(j)\). The finite sum over \(i+j=q\) gives the stated bound. Vertical equivalence follows from the two contractions used above. ARCH-136-TW-TENSOR transfers the horizontal assertion to \(\mathcal T(C)\).
\end{proof}

\section{A limitation on saturation inside standard rational forms}

\paragraph{ARCH-136-INTEGRAL-PERIOD-MAXIMALITY.}
Let \(B_{\mathbb Z}\subset E\subset B_{\mathbb Q}\) be a simplicial bigraded differential \(S=\Gamma_{\mathbb Z}(T)\)-subalgebra. Suppose that for every \(r\) there is \(W_r\) such that, at every simplex and weight \(w\geq W_r\), the coefficient map identifies \(H^0(E^{w,*})\) with \(\mathbb ZT^{[w]}\), and \(H^j(E^{w,*})=0\) for \(1\leq j\leq r\). Then \(E=B_{\mathbb Z}\). No horizontal admissibility condition is required.

\begin{proof}
Use rational simplex integration, normalized by
\[
 \int_{\Delta^n}
 \gamma_{a_0}(x_0)\cdots\gamma_{a_n}(x_n)\,
 dx_1\cdots dx_n=T^{[|a|+n]}.
\]
This is the usual polynomial integral with \(x_0=T-\sum_{i\geq1}x_i\); the beta-integral formula gives the displayed identity, and Stokes holds.

Choose \(W_n\) to dominate the vertical thresholds through \(n\). If \(w\geq W_n\) and \(\alpha\in E_n^{w,n}\), then \(\alpha\) is closed because the ambient algebra has no forms of degree greater than \(n\). For \(n>0\), write \(\alpha=d\beta\) in \(E_n\). Each face of \(\beta\) is a top form on an \((n-1)\)-simplex. Induction on \(n\), followed by Stokes, gives
\(\int\alpha\in\mathbb ZT^{[w]}\). For \(n=0\), this is precisely the assumed coefficient statement, since \(E_0\) has no positive form degrees.

The same integrality holds at every weight. Otherwise suppose
\(\int\alpha=cT^{[L]}\), with \(c\in\mathbb Q\setminus\mathbb Z\), and choose a prime \(p\) with \(v_p(c)<0\). For \(p^a>L\),
\[
 \int T^{[p^a]}\alpha
 =\binom{L+p^a}{L}c\,T^{[L+p^a]}.
\]
The binomial coefficient is \(1\) modulo \(p\), so this remains nonintegral at arbitrarily large weights, contradicting the preceding paragraph.

Fix \(f\in E_n^{w,q}\) and expand it in the rational barycentric divided-power basis:
\[
 f=\sum_{\substack{|I|=q\\|a|=w-q}}
 c_{I,a}\,\gamma_{a_0}(x_0)\cdots\gamma_{a_n}(x_n)\,dx_I .
\]
For fixed \(I\), multiply by
\(\gamma_{b_0}(x_0)\cdots\gamma_{b_n}(x_n)\,dx_{I^c}\in B_{\mathbb Z}\).
All other exterior summands vanish. Up to a fixed orientation sign, the resulting integral has coefficient
\[
 P_I(b)=\sum_{|a|=w-q}c_{I,a}
          \prod_{j=0}^n\binom{a_j+b_j}{a_j}.
\]
This is an integer for every \(b\in\mathbb N^{n+1}\). Therefore every iterated forward difference at \(0\) is an integer. The top-order difference
\(\Delta_0^{a_0}\cdots\Delta_n^{a_n}P_I(0)\) equals \(c_{I,a}\): Vandermonde expands each factor in the Newton basis \(\binom{b_j}{c_j}\), and a summand indexed by \(a'\) with \(|a'|=|a|\) contributes to multiindex \(a\) only if \(a'_j\geq a_j\) for all \(j\), which forces \(a'=a\). Thus all coefficients are integers and \(f\in B_{\mathbb Z}\).
\end{proof}

\paragraph{Scope.}
The last theorem rules out a proper admissible enlargement of standard integral forms inside standard rational forms. It does not apply to an enlargement inside \(\mathcal T(C)_{\mathbb Q}\), which has forms above the original simplex dimension; its top-dimensional coefficient need not be closed. Nor does it imply that the image of an admissible theory under a map to \(B_{\mathbb Q}\) is admissible.

